\section{Minimum Insertions to Make a String Palindrome}
\label{sec:dp-min-insertions-palindrome}

\problemstatement{Given a string $s$ of length $n$, find the minimum
number of characters that must be \textbf{inserted} (anywhere in the
string) to make it a palindrome.}

\begin{patternbox}
A second reduction layered directly on top of
Section~\ref{sec:dp-longest-palindromic-subsequence}: the characters
\emph{not} part of the longest palindromic subsequence are exactly the
ones that need a ``partner'' inserted to balance them into a palindrome,
so the answer is simply $n$ minus the LPS length.
\end{patternbox}

\subsection*{Understanding the reduction}

Any palindromic subsequence of $s$, of length $k$, can be extended to a
full palindrome covering all of $s$ by inserting, for every character of
$s$ \emph{not} part of that subsequence, a matching mirrored character
on the opposite side. The longest palindromic subsequence minimizes the
number of characters left over needing such a partner, so:
\[
  \textit{minInsertions} = n - |\textit{LPS}(s)|.
\]

\subsection*{All Four Approaches}

Every stage reuses Section~\ref{sec:dp-longest-palindromic-subsequence}'s
implementation (itself a reduction to LCS) unchanged.

\begin{lstlisting}
#include <bits/stdc++.h>
using namespace std;

int lcsSpaceOptimized(string& s1, string& s2) {
    int n = s1.size(), m = s2.size();
    vector<int> prev(m + 1, 0), current(m + 1, 0);

    for (int i = 1; i <= n; i++) {
        for (int j = 1; j <= m; j++) {
            if (s1[i - 1] == s2[j - 1]) {
                current[j] = 1 + prev[j - 1];
            } else {
                current[j] = max(prev[j], current[j - 1]);
            }
        }
        prev = current;
    }
    return prev[m];
}

int minInsertionsForPalindrome(string s) {
    string r = s;
    reverse(r.begin(), r.end());
    int lpsLen = lcsSpaceOptimized(s, r);
    return (int)s.size() - lpsLen;
}

int main() {
    cout << minInsertionsForPalindrome("zzazz") << endl;  // 0 (already a palindrome)
    cout << minInsertionsForPalindrome("mbadm") << endl;  // 2
    return 0;
}
\end{lstlisting}

\subsection*{Dry run}

\dryrun{Take $s = \texttt{"mbadm"}$ ($n=5$).}

The longest palindromic subsequence of \texttt{"mbadm"} is
\texttt{"mam"} (length $3$): the characters \texttt{m}, \texttt{a},
\texttt{m} appear in this order at indices $0,2,4$, and no length-$4$
palindromic subsequence exists (no other character repeats).
$\textit{minInsertions} = 5-3=2$.

\begin{complexitybox}
\textbf{Time Complexity.} $O(n^2)$, inherited from the LCS-based LPS
computation.
\textbf{Space Complexity.} $O(n)$ for the rolling-row LCS.
\end{complexitybox}

\begin{takeawaybox}
This problem closes a three-link reduction chain that is worth seeing
as a whole: Minimum Insertions for Palindrome reduces to Longest
Palindromic Subsequence, which reduces to Longest Common Subsequence.
Recognizing chains of reductions like this -- rather than treating each
problem as requiring its own fresh derivation -- is the single most
efficient way to work through a sheet of dozens of nominally distinct DP
problems.
\end{takeawaybox}
